% ─────────────────────────────────────────────────────────────────────────────
% LabBank — ALM System Presentation
% Compile: pdflatex labbank_presentation.tex  (twice for ToC)
% Assets:  run make_beamer_assets.py first → produces beamer_assets/*.png
% ─────────────────────────────────────────────────────────────────────────────
\documentclass[aspectratio=169,10pt]{beamer}

\usetheme{Madrid}
\usecolortheme{default}

% ── Custom colours ────────────────────────────────────────────────────────────
\definecolor{labblue}{RGB}{21,  101, 192}   % asset blue
\definecolor{labred} {RGB}{198,  40,  40}   % liability red
\definecolor{labgrn} {RGB}{ 46, 125,  50}   % net green
\definecolor{labpur} {RGB}{106,  27, 154}   % equity purple
\definecolor{labdark}{RGB}{ 26,  35, 126}   % deep navy (header)
\definecolor{labgrey}{RGB}{ 96, 125, 139}   % muted grey

\setbeamercolor{palette primary}  {bg=labdark, fg=white}
\setbeamercolor{palette secondary}{bg=labblue, fg=white}
\setbeamercolor{palette tertiary} {bg=labblue!60, fg=white}
\setbeamercolor{structure}        {fg=labblue}
\setbeamercolor{frametitle}       {bg=labdark, fg=white}
\setbeamercolor{title}            {bg=labdark, fg=white}
\setbeamercolor{alerted text}     {fg=labred}
\setbeamercolor{block title}      {bg=labblue!15, fg=labdark}
\setbeamercolor{block body}       {bg=labblue!5,  fg=black}
\setbeamercolor{block title alerted}{bg=labred!15, fg=labred!80!black}
\setbeamercolor{block body alerted} {bg=labred!5,  fg=black}
\setbeamercolor{block title example}{bg=labgrn!15, fg=labgrn!70!black}
\setbeamercolor{block body example} {bg=labgrn!5,  fg=black}

% ── Packages ─────────────────────────────────────────────────────────────────
\usepackage{graphicx}
\usepackage{booktabs}
\usepackage{amsmath}
\usepackage{xcolor}
\usepackage{tikz}
\usetikzlibrary{arrows.meta, positioning, shapes.geometric}
\usepackage{multirow}
\usepackage{array}
\usepackage{microtype}
\usepackage{hyperref}

\hypersetup{colorlinks=false}
\setbeamertemplate{navigation symbols}{}
\setbeamertemplate{footline}[frame number]

% ── Convenience macros ────────────────────────────────────────────────────────
\newcommand{\imgfull}[2][0.96\textwidth]{%
  \begin{center}\includegraphics[width=#1]{beamer_assets/#2}\end{center}}
\newcommand{\pos}[1]{\textcolor{labgrn}{\textbf{#1}}}
\newcommand{\neg}[1]{\textcolor{labred}{\textbf{#1}}}
\newcommand{\hi} [1]{\textcolor{labblue}{\textbf{#1}}}
\newcommand{\formula}[1]{\colorbox{labblue!8}{\(#1\)}}

% ─────────────────────────────────────────────────────────────────────────────
% TITLE
% ─────────────────────────────────────────────────────────────────────────────
\title[LabBank — ALM System]{\textbf{LabBank}\\[4pt]
  \large A Laboratory Banking ALM System}
\subtitle{Balance Sheet · Cash Flows · IRRBB · Liquidity Risk · Reporting}
\author{ALM \& Risk Analytics}
\date{Report date: 2024-12-31}
\institute{LabBank Project}

% ─────────────────────────────────────────────────────────────────────────────
\begin{document}
% ─────────────────────────────────────────────────────────────────────────────

% ── Slide 1 — Title ───────────────────────────────────────────────────────────
\begin{frame}[plain]
  \titlepage
  \vspace{-6pt}
  \begin{center}
    \small\textcolor{labgrey}{Currency: PLN \quad|\quad
      Risk framework: EBA IRRBB Guidelines (2018)\quad|\quad
      Audience: ALM / Risk Analytics}
  \end{center}
\end{frame}

% ── Slide 2 — Outline ─────────────────────────────────────────────────────────
\begin{frame}{Agenda}
  \tableofcontents
\end{frame}

% =============================================================================
\section{Project Overview}
% =============================================================================

% ── Slide 3 — What is LabBank ─────────────────────────────────────────────────
\begin{frame}{What is LabBank?}

  \begin{columns}[T]
    \column{0.54\textwidth}
      \textbf{LabBank is an end-to-end ALM modelling system} built around a
      synthetic but realistic bank balance sheet.
      It computes the key regulatory and management risk metrics from first
      principles — starting from raw transactions and market curves, producing
      auditable, SQL-backed results.

      \vspace{8pt}
      \textbf{Why it exists}
      \begin{itemize}
        \item Demonstrate how IRRBB and liquidity risk measures can be derived
              from elementary cash flow aggregation.
        \item Provide a \textit{transparent}, fully reproducible reference
              implementation — every number traces back to a SQL table or a
              Python calculation.
        \item Support education in quantitative ALM for practitioners who want
              to understand the mechanics rather than treat metrics as black-box
              outputs.
      \end{itemize}

    \column{0.42\textwidth}
      \begin{block}{What LabBank computes}
        \small
        \begin{tabular}{ll}
          \toprule
          \textbf{Module} & \textbf{Output} \\
          \midrule
          Balance sheet   & Transactions, rates \\
          Cash flows      & CF schedules, NII \\
          IR derivatives  & Swap CFs, gap adj. \\
          IRRBB           & NII, EVE, EBA SOT \\
          Liquidity       & LCR, NSFR \\
          Reporting       & HTML, PDF, Power BI \\
          \bottomrule
        \end{tabular}
      \end{block}

      \begin{alertblock}{Scope note}
        \small Balance sheet \textbf{optimisation} is the planned next phase.
        Methods are not yet fully implemented and are \textit{outside} the
        scope of this presentation.
      \end{alertblock}
  \end{columns}
\end{frame}

% =============================================================================
\section{System Architecture}
% =============================================================================

% ── Slide 4 — Architecture ────────────────────────────────────────────────────
\begin{frame}{System Architecture}

  \imgfull{pipeline.png}

  \vspace{2pt}
  \small
  Each module writes to a dedicated \textbf{SQL Server schema} and exposes
  clean tabular outputs to downstream modules.
  The final reporting layer consumes all schemas to produce HTML/PDF reports.
  SQL acts as the integration bus: modules are independently runnable and
  results are fully auditable at every step.
\end{frame}

% ── Slide 5 — SQL Schema ──────────────────────────────────────────────────────
\begin{frame}{Data Model — SQL Schemas}

  \imgfull[0.88\textwidth]{sql_schema.png}

  \small\vspace{2pt}
  Six schemas separate concerns cleanly.
  \texttt{dbo} and \texttt{schemat} hold raw inputs; \texttt{mkt} holds
  market curves; \texttt{cf}, \texttt{irrbb}, and \texttt{results} hold
  calculated outputs.
  Every calculation reads from upstream schemas — no hidden side-effects.
\end{frame}

% =============================================================================
\section{Balance Sheet}
% =============================================================================

% ── Slide 6 — BS Construction ─────────────────────────────────────────────────
\begin{frame}{Balance Sheet Construction}

  \begin{columns}[T]
    \column{0.48\textwidth}
      \textbf{Input} (\texttt{bank\_data.xlsx})
      \begin{itemize}
        \item Product definitions: code, name, side (A/L).
        \item Balance amounts by product and currency.
        \item Contractual rates, maturity dates, repricing frequencies.
        \item Behavioural parameters: prepayment speeds for loans,
              decay rates for non-maturity deposits (NMD).
        \item Equity (Tier 1 capital) — used as regulatory denominator
              in the SOT.
      \end{itemize}

      \vspace{6pt}
      \textbf{Rate types}
      \begin{description}[labelwidth=1.8cm]
        \item[\texttt{Fixed (F)}]  Rate contractually locked until maturity.
        \item[\texttt{Variable (V)}] Resets periodically to WIBOR + spread.
        \item[\texttt{Admin (A)}]  Floored at 0\%; moves with policy at the
                                    bank's discretion (e.g.\ current accounts).
      \end{description}

    \column{0.48\textwidth}
      \textbf{Output tables produced}

      \begin{exampleblock}{\texttt{dbo.transactions}}
        Master balance sheet — one row per product per report date.
        Columns: \texttt{balance\_amt}, \texttt{client\_rt},
        \texttt{rate\_type}, \texttt{bs\_side}, \texttt{currency}.
      \end{exampleblock}

      \begin{exampleblock}{\texttt{schemat.loans / deposits}}
        Product-level schedules with individual contractual maturities,
        coupon dates, and behavioural run-off parameters.
      \end{exampleblock}

      \vspace{4pt}
      \textbf{Multi-currency:} PLN is the reporting currency.
        EUR and USD positions are modelled separately with their own
        discount and forward curves.
  \end{columns}
\end{frame}

% ── Slide 7 — BS Structure chart ──────────────────────────────────────────────
\begin{frame}{Balance Sheet Structure}

  \imgfull{bs_structure.png}

  \vspace{2pt}
  \small
  \begin{columns}
    \column{0.48\textwidth}
      \textbf{Asset side:} mortgages (fixed \& float), consumer loans,
      bond/securities portfolio, interbank placements.
    \column{0.48\textwidth}
      \textbf{Liability side:} term deposits, savings accounts, current
      accounts, issued bonds, plus equity (Tier 1 capital).
  \end{columns}
  \vspace{4pt}
  Key metrics: \hi{NIM 5.35\%} \quad|\quad
               Avg asset yield \hi{8.16\%} \quad|\quad
               Cost of funds \hi{3.35\%} \quad|\quad
               Interest spread \hi{4.81\%}
\end{frame}

% ── Slide 8 — Rate type exposure ──────────────────────────────────────────────
\begin{frame}{Rate Type Exposure}

  \imgfull{rate_type.png}

  \small\vspace{2pt}
  Rate type determines \textbf{how fast} a position reprices when market
  rates move:
  \begin{itemize}
    \item \textbf{Fixed-rate positions} produce stable income regardless of
          rate moves — but their \textit{economic value} changes (EVE risk).
    \item \textbf{Variable-rate positions} reprice quickly — NII changes
          almost immediately when WIBOR shifts (NII risk).
    \item \textbf{Admin-rate liabilities} (current accounts) are floored at
          0\% and move slowly, creating a structural hedge in falling-rate
          environments.
  \end{itemize}
  The mix between fixed and variable drives the bank's position on the
  \textbf{NII vs EVE trade-off}, central to IRRBB management.
\end{frame}

% =============================================================================
\section{Cash Flows}
% =============================================================================

% ── Slide 9 — CF Engine ───────────────────────────────────────────────────────
\begin{frame}{Cash Flow Engine}

  \begin{columns}[T]
    \column{0.50\textwidth}
      \textbf{What the CF engine does}

      For each product on the balance sheet it generates a full schedule of:
      \begin{enumerate}
        \item \textbf{Capital cash flows} — contractual amortisation plus
              behavioural adjustments (early repayments, deposit decay).
        \item \textbf{Interest cash flows} — coupon payments at the
              contractual client rate, separately tracked by tenor bucket.
      \end{enumerate}

      \vspace{6pt}
      \textbf{Behavioural modelling}
      \begin{itemize}
        \item \textit{Loans:} constant prepayment rate (CPR) applied monthly
              to the outstanding balance.
        \item \textit{Non-maturity deposits (NMD):} exponential decay model
              — a fraction of the balance is assumed to leave each month,
              the remainder is treated as sticky long-term funding.
        \item \textit{Admin-rate floor:} interest on current accounts is
              capped at 0\% and never passes through to the client as
              a negative rate.
      \end{itemize}

    \column{0.46\textwidth}
      \textbf{Output: \texttt{cf.nii\_base\_scenario}}

      One row per (product, tenor bucket, scenario).  Key columns:

      \vspace{4pt}
      \small
      \begin{tabular}{ll}
        \toprule
        Column & Meaning \\
        \midrule
        \texttt{beh\_outstanding} & Behavioural balance outstanding \\
        \texttt{client\_rt}       & Contracted client rate \\
        \texttt{fwd\_rt}          & Market forward rate (same tenor) \\
        \texttt{margin}           & \texttt{client\_rt} $-$ \texttt{fwd\_rt} \\
        \texttt{nii\_interest}    & Income from existing CFs (locked) \\
        \texttt{nii\_renewal}     & Income from maturing capital \\
                                  & reinvested at market rates \\
        \texttt{nii\_total}       & Sum of the two \\
        \bottomrule
      \end{tabular}

      \vspace{6pt}
      \normalsize
      This table is the \textit{single source of truth} for both NII
      reporting and scenario analysis.
  \end{columns}
\end{frame}

% ── Slide 10 — NII Bridge ─────────────────────────────────────────────────────
\begin{frame}{NII: Locked Interest vs Renewal}

  \imgfull{nii_bridge.png}

  \small\vspace{2pt}
  \begin{columns}
    \column{0.48\textwidth}
      \textbf{Locked interest} (\texttt{nii\_interest}): income from cash
      flows already contracted — rate and notional are fixed.
      This part is \textit{immune to rate shocks}.
    \column{0.48\textwidth}
      \textbf{Renewal} (\texttt{nii\_renewal}): income from maturing capital
      that is reinvested at current market rates within the 1-year horizon.
      This part \textit{reprices with the curve} and drives NII sensitivity.
  \end{columns}
\end{frame}

% =============================================================================
\section{Repricing Gap \& IR Derivatives}
% =============================================================================

% ── Slide 11 — Repricing gap ──────────────────────────────────────────────────
\begin{frame}{Repricing Gap and Interest Rate Swaps}

  \imgfull{repricing_gap.png}

  \small\vspace{2pt}
  \textbf{Repricing gap} = net volume of assets minus liabilities
  repricing in each tenor bucket.
  A \pos{positive} gap means more assets than liabilities reprice in that
  bucket — the bank \textit{benefits} from a rate increase.
  A \neg{negative} gap signals liability-sensitive exposure — the bank is
  hurt by rising rates in that bucket.

  \vspace{4pt}
  \textbf{IRS book:} interest rate swaps shift the gap without changing
  the balance sheet.
  A \textit{receive-fixed} swap converts variable-rate liabilities into
  fixed-rate obligations, reducing short-term asset sensitivity.
\end{frame}

% =============================================================================
\section{IRRBB — NII and EVE}
% =============================================================================

% ── Slide 12 — NII Waterfall ──────────────────────────────────────────────────
\begin{frame}{Net Interest Income — Product Breakdown}

  \imgfull{nii_waterfall.png}

  \small\vspace{2pt}
  Interest income accumulates from asset products (mortgages, loans,
  securities), then interest expense from liability products (deposits,
  bonds) is subtracted to arrive at Net Interest Income.
  This waterfall makes the \textbf{product-level contribution} explicit —
  key for pricing decisions and product strategy.
\end{frame}

% ── Slide 13 — EVE ────────────────────────────────────────────────────────────
\begin{frame}{Economic Value of Equity (EVE)}

  \begin{columns}[T]
    \column{0.46\textwidth}
      \textbf{Definition}

      EVE is the present value of all future cash flows under a
      \textit{run-off} assumption — no new business, existing book
      winds down to maturity:
      \[
        \mathrm{EVE} = \sum_{i} \mathrm{CF}_i \cdot d(t_i)
      \]
      where $d(t_i)$ is the discount factor at tenor $t_i$.

      \vspace{8pt}
      \textbf{What ΔEVE tells us}

      When the yield curve shifts, the discount factors change and so
      does the present value of the book.
      \[
        \Delta\mathrm{EVE} = \mathrm{EVE}_{\text{shocked}}
                           - \mathrm{EVE}_{\text{base}}
      \]
      A \neg{negative} ΔEVE means the bank's economic net worth falls
      under that rate scenario — long-maturity fixed assets lose value
      faster than liabilities shorten.

    \column{0.50\textwidth}
      \imgfull[0.97\textwidth]{eve_results.png}
  \end{columns}
\end{frame}

% =============================================================================
\section{EBA Shock Scenarios \& SOT}
% =============================================================================

% ── Slide 14 — Shock curves ───────────────────────────────────────────────────
\begin{frame}{EBA IRRBB Shock Scenarios}

  \imgfull{shock_curves.png}

  \vspace{2pt}
  \small
  EBA Guidelines (2018) prescribe \textbf{six supervisory shock scenarios}
  applied to each currency's yield curve:

  \vspace{4pt}
  \begin{columns}
    \column{0.48\textwidth}
      \begin{itemize}
        \item \hi{Parallel up} (+250 bps across all tenors)
        \item \hi{Parallel down} (−250 bps across all tenors)
        \item \hi{Steepener} (short rates down, long rates up)
      \end{itemize}
    \column{0.48\textwidth}
      \begin{itemize}
        \item \hi{Flattener} (short rates up, long rates down)
        \item \hi{Short rate up} (+350 bps at the short end)
        \item \hi{Short rate down} (−350 bps at the short end)
      \end{itemize}
  \end{columns}

  \vspace{4pt}
  Shocks are floored at 0\% for each tenor.
  Each scenario produces a full set of ΔEVE and ΔNII figures
  used in the Supervisory Outlier Test.
\end{frame}

% ── Slide 15 — SOT ────────────────────────────────────────────────────────────
\begin{frame}{Supervisory Outlier Test (SOT)}

  \begin{columns}[T]
    \column{0.44\textwidth}
      \textbf{Methodology}

      The SOT applies a \textbf{50\% haircut per tenor bucket} to any
      scenario leg that produces a \textit{gain}:
      \[
        \Delta\mathrm{EVE}^{\text{reg}} = \sum_{b}
          \min\!\bigl(\Delta\mathrm{EVE}_b,\;0\bigr)
          + 0.5 \cdot \max\!\bigl(\Delta\mathrm{EVE}_b,\;0\bigr)
      \]
      The haircut is applied \textit{at bucket level}, not on the total —
      preventing netting of gains in one bucket against losses in another.

      \vspace{6pt}
      \textbf{Thresholds (EBA)}
      \begin{itemize}
        \item EVE SOT: $\Delta\mathrm{EVE}^{\text{reg}} / T_1 < -15\%$
              \neg{\(\Rightarrow\) BREACH}
        \item NII SOT: $\Delta\mathrm{NII}^{\text{reg}} / T_1 < -5\%$
              \neg{\(\Rightarrow\) BREACH}
      \end{itemize}
      Results are written to \texttt{irrbb.irrbb\_report}.

    \column{0.52\textwidth}
      \imgfull[0.97\textwidth]{eba_sot.png}
  \end{columns}
\end{frame}

% =============================================================================
\section{Liquidity Risk}
% =============================================================================

% ── Slide 16 — LCR / NSFR ────────────────────────────────────────────────────
\begin{frame}{Liquidity Coverage Ratio \& Net Stable Funding Ratio}

  \begin{columns}[T]
    \column{0.44\textwidth}
      \textbf{LCR — 30-day stress horizon}
      \[
        \mathrm{LCR} = \frac{\mathrm{HQLA}}
                            {\text{Net outflows}_{30\text{d}}}
                       \geq 100\%
      \]
      HQLA (High-Quality Liquid Assets) = Level 1 liquid assets:
      cash, central bank reserves, sovereign bonds.
      Net outflows = stressed deposit run-off and interbank maturities
      within 30 days.

      \vspace{8pt}
      \textbf{NSFR — 1-year structural funding}
      \[
        \mathrm{NSFR} = \frac{\mathrm{ASF}}{\mathrm{RSF}} \geq 100\%
      \]
      ASF (Available Stable Funding) is weighted by funding stability
      (equity \& long-term debt score highest).
      RSF (Required Stable Funding) is weighted by asset liquidity
      (loans and illiquid assets score highest).

    \column{0.52\textwidth}
      \imgfull[0.97\textwidth]{lcr_nsfr.png}
  \end{columns}
\end{frame}

% ── Slide 17 — Liquidity Gap ──────────────────────────────────────────────────
\begin{frame}{Behavioural Liquidity Gap}

  \imgfull{liq_gap.png}

  \small\vspace{2pt}
  The liquidity gap shows the net cash position across the full
  tenor ladder up to 60 months, using \textbf{behavioural} rather than
  contractual cash flows (loan prepayments included, NMD decay applied).
  A sustained \neg{negative cumulative gap} identifies the point at which
  the bank would need to seek additional funding — the primary input for
  contingency funding planning.
\end{frame}

% =============================================================================
\section{Reporting Pipeline}
% =============================================================================

% ── Slide 18 — Reporting ──────────────────────────────────────────────────────
\begin{frame}{Integrated Reporting}

  \begin{columns}[T]
    \column{0.48\textwidth}
      \textbf{bank\_report (ALM \& Compliance)}

      Jupyter notebook executed against live SQL data and
      exported to \texttt{HTML / PDF}:
      \begin{itemize}
        \item Yield curves and EBA shocked scenarios
        \item Repricing gap ladder with IRS overlay
        \item Behavioural liquidity gap (60M horizon)
        \item NII by scenario (all 7 shocks)
        \item EVE by scenario (run-off PV)
        \item EBA SOT: ΔEVE \% and ΔNII \% vs Tier 1
        \item LCR \& NSFR by currency
        \item Non-linearity analysis (convexity check)
      \end{itemize}

    \column{0.48\textwidth}
      \textbf{finance\_report (Finance / CFO)}

      Programmatically generated notebook via
      \texttt{make\_finance\_report.py}:
      \begin{itemize}
        \item Executive P\&L summary (NIM, spread, NII)
        \item Interest income \& expense waterfall
        \item Product margin vs benchmark (bps spread)
        \item Funding cost structure (pie + cost rates)
        \item Fixed vs variable rate mix
        \item NII bridge: locked interest vs renewal
        \item Margin profile over time (by tenor bucket)
        \item Rate sensitivity (business view)
      \end{itemize}

      \vspace{4pt}
      \textbf{Power BI:} \texttt{ir\_gap.pbix} provides interactive
      repricing gap exploration connected to the live SQL database.
  \end{columns}
\end{frame}

% =============================================================================
\section{NII Scenario Analysis}
% =============================================================================

% ── Slide 19 — NII scenarios ──────────────────────────────────────────────────
\begin{frame}{NII Across EBA Shock Scenarios}

  \imgfull{nii_scenarios.png}

  \small\vspace{2pt}
  NII decomposes into \textbf{locked interest} (immune to shocks, shown
  solid) and \textbf{renewal} (rate-sensitive, shown hatched).
  The delta shown above each bar is the change versus the base scenario.
  \textbf{Asset-sensitive balance sheets} benefit from parallel-up shocks;
  \textbf{liability-sensitive} balance sheets benefit from parallel-down.
\end{frame}

% =============================================================================
\section{Conclusions \& Roadmap}
% =============================================================================

% ── Slide 20 — Conclusions ────────────────────────────────────────────────────
\begin{frame}{Conclusions \& What Comes Next}

  \begin{columns}[T]
    \column{0.52\textwidth}
      \textbf{What LabBank delivers today}
      \begin{itemize}
        \item Full end-to-end ALM pipeline from raw transactions to
              regulatory risk metrics — reproducible, SQL-backed,
              auditable at every step.
        \item EBA-compliant IRRBB: correct shock construction, per-bucket
              SOT haircut, multi-currency curves.
        \item Behavioural realism: loan prepayment, NMD decay,
              admin-rate floor — not purely contractual.
        \item NII and EVE computed consistently from the \textit{same}
              cash flow schedules, eliminating reconciliation gaps.
        \item Dual reporting: ALM/Compliance view and Finance/CFO view,
              both auto-generated from live SQL.
        \item Interactive gap analysis via Power BI.
      \end{itemize}

    \column{0.44\textwidth}
      \textbf{Planned: balance sheet optimisation}

      The next phase will add a fast approximation layer
      (\texttt{optimize\_prep}) that pre-computes vectorised
      sensitivity formulas, enabling:
      \begin{itemize}
        \item Portfolio optimisation subject to SOT, LCR, and NSFR
              constraints.
        \item Efficient frontier analysis for NIM vs ΔEVE.
        \item Hedging strategy evaluation (IRS overlay selection).
      \end{itemize}

      Methods are currently in development and not yet
      production-ready.

      \vspace{8pt}
      \begin{exampleblock}{LabBank as an educational tool}
        \small Every formula in this presentation has a direct SQL or
        Python equivalent in the codebase.
        The system is designed to make ALM maths \textit{concrete} —
        not abstract — by connecting equations to actual data rows.
      \end{exampleblock}
  \end{columns}
\end{frame}

% ── Final slide ───────────────────────────────────────────────────────────────
\begin{frame}[plain]
  \begin{center}
    \vspace{20pt}
    {\Large\textbf{\textcolor{labdark}{LabBank}}}\\[8pt]
    {\large\textcolor{labgrey}{Laboratory Banking ALM System}}\\[20pt]
    \begin{tabular}{ll}
      \textbf{Modules:}    & Balance sheet · Cash flows · IRRBB · Liquidity \\[4pt]
      \textbf{Framework:}  & EBA IRRBB Guidelines (2018) \\[4pt]
      \textbf{Technology:} & Python · SQL Server · Jupyter · Power BI \\[4pt]
      \textbf{Next:}       & Balance sheet optimisation \\
    \end{tabular}
    \vspace{20pt}\\
    \textcolor{labgrey}{\small Report date: 2024-12-31 \quad|\quad Currency: PLN}
  \end{center}
\end{frame}

\end{document}
